\documentclass[12pt,a4paper]{article}

% ── Fonts (LuaLaTeX) ──────────────────────────────────────────────────────────
\usepackage{fontspec}
\usepackage{unicode-math}
\setmainfont{Linux Libertine O}
\setmathfont{TeX Gyre Pagella Math}
\newfontfamily\shavfont{Linux Libertine O}
\newfontfamily\igprimfont{Everson Mono}
\setmonofont[Scale=0.85]{Everson Mono}

% ── Layout ────────────────────────────────────────────────────────────────────
\usepackage[top=1in, bottom=1in, left=1.1in, right=1.1in, headheight=15pt]{geometry}
\usepackage{microtype}
\usepackage{parskip}
\usepackage[english]{babel}

% ── Headers ───────────────────────────────────────────────────────────────────
\usepackage{fancyhdr}
\pagestyle{fancy}
\fancyhf{}
\fancyhead[L]{\leftmark}
\fancyhead[R]{\thepage}
\renewcommand{\headrulewidth}{0.4pt}

% ── Section styling ───────────────────────────────────────────────────────────
\usepackage{titlesec}
\titleformat{\section}{\Large\bfseries}{\thesection}{1em}{}
\titleformat{\subsection}{\large\bfseries}{\thesubsection}{1em}{}

% ── Lists ─────────────────────────────────────────────────────────────────────
\usepackage{enumitem}
\setlist[itemize]{topsep=0.5ex, itemsep=0.25ex, parsep=0pt}
\setlist[enumerate]{topsep=0.5ex, itemsep=0.25ex, parsep=0pt}

% ── Math and tables ───────────────────────────────────────────────────────────
\usepackage{amsmath,amsthm}
\usepackage{mathtools}
\usepackage{booktabs}
\usepackage{array}
\usepackage{tabularx}

% ── Code listings ─────────────────────────────────────────────────────────────
\usepackage{listings}
\usepackage{xcolor}
\definecolor{codebg}{RGB}{248,248,248}
\definecolor{codeframe}{RGB}{200,200,200}
\definecolor{codekeyword}{RGB}{0,0,180}
\definecolor{codecomment}{RGB}{0,128,0}
\lstset{
  basicstyle=\ttfamily\footnotesize,
  backgroundcolor=\color{codebg},
  frame=single,
  rulecolor=\color{codeframe},
  keywordstyle=\color{codekeyword}\bfseries,
  commentstyle=\color{codecomment}\itshape,
  breaklines=true,
  captionpos=b,
  morekeywords={axiom,def,theorem,noncomputable,structure,where,namespace,
                open,import,section,end,fun,exact,intro,simp,rfl,decide,
                constructor,obtain,refine,cases,by,have,exact,use,rw},
}

% ── Cross-references ──────────────────────────────────────────────────────────
\usepackage{hyperref}
\usepackage{cleveref}
\hypersetup{colorlinks=true, linkcolor=blue, urlcolor=cyan, citecolor=blue}

% ── Theorem environments ──────────────────────────────────────────────────────
\newtheorem{theorem}{Theorem}[section]
\newtheorem{lemma}[theorem]{Lemma}
\newtheorem{proposition}[theorem]{Proposition}
\newtheorem{corollary}[theorem]{Corollary}
\newtheorem{definition}[theorem]{Definition}
\newtheorem{remark}[theorem]{Remark}

% ── Custom notation ───────────────────────────────────────────────────────────
\newcommand{\BFour}{\mathbf{B}_4}
\newcommand{\WH}{\mathrm{WH}}
\newcommand{\SIC}{\mathrm{SIC}}
\newcommand{\frob}{\mu \circ \delta = \mathrm{id}}
\newcommand{\Lean}{\textsc{Lean\,4}}
\newcommand{\frobInner}{\langle\,\cdot\,,\cdot\,\rangle_{\!\mathbb{N}}}

% ─────────────────────────────────────────────────────────────────────────────

\begin{document}

\title{\textbf{The Symmetric Informationally-Complete Measurement is
Natively Four-Valued:\\
A Machine-Verified Belnap-Multilattice Construction for all $d = 2^n$}}

\author{C.\ L.\ Mills (Lando\,$\otimes$\,$\odot$-operator)%
\thanks{The author thanks Harry T.\ Larson \cite{Larson1961}.}}

\date{July 2026}

\maketitle

\begin{abstract}
A symmetric informationally-complete positive operator-valued measure
(SIC-POVM) is the optimally informative quantum measurement, and whether one
exists in every finite dimension is the content of Zauner's conjecture, open
since 1999. We prove, with a fully machine-checked \Lean{} development
carrying zero axioms and zero \texttt{sorry}s, that the complete structural
content of a Weyl--Heisenberg covariant SIC-POVM is already present, in every
dimension $d = 2^n$, in the four-valued paraconsistent Belnap lattice: the
Weyl--Heisenberg orbit of the all-\emph{Both} fiducial has exactly $4^n = d^2$
elements, satisfies Frobenius closure $\frob$, and is equiangular with respect
to an integer-valued Frobenius inner product. This structural SIC-POVM is
unconditional. We then isolate exactly what remains: the passage from this
structure to a fiducial \emph{vector} in $\mathbb{C}^d$ with constant unit
overlap is a representation problem, and that problem is Zauner's conjecture
verbatim. The openness therefore attaches to the complex-Hilbert
representation, not to the SIC-POVM itself; the arithmetic-geometry route to
the vector, via the mixed-signature Stark conjecture for the ray class fields
$\mathbb{Q}(\sqrt{d(d-2)})$, is presented as one such representation and its
irreducible inputs are named. The upshot is a structural inversion: the ideal
quantum measurement is not a fact about complex Hilbert space that four-valued
logic happens to model, but a fact about four-valued logic of which complex
Hilbert space is one representation.
\end{abstract}

\tableofcontents
\newpage

%──────────────────────────────────────────────────────────────────────────────
\section{Introduction}
%──────────────────────────────────────────────────────────────────────────────

Fix a finite dimension $d$. A rank-one SIC-POVM is a set of $d^2$ unit vectors
$\{\psi_j\}$ in $\mathbb{C}^d$ whose pairwise overlaps are constant,
$|\langle \psi_j, \psi_k\rangle|^2 = \tfrac{1}{d+1}$ for $j \neq k$. Such a
configuration is the tightest possible informationally-complete measurement,
and it is the natural resolution of the identity used throughout quantum
tomography and quantum Bayesianism. Zauner's conjecture \cite{Zauner1999}
asserts that a Weyl--Heisenberg covariant SIC-POVM exists in every dimension.
It is supported by exact solutions in many dimensions and high-precision
numerical solutions in many more, yet a general existence proof has resisted
since 1999.

This paper does not resolve Zauner's conjecture in $\mathbb{C}^d$. It does
something that, for the dimensions $d = 2^n$, we argue is prior to it: it shows
that the entire structural skeleton a SIC-POVM is asked to instantiate, its
symmetry, its cardinality, its equiangularity, and its self-duality, exists
unconditionally and can be exhibited by hand, in a four-valued logic, with no
appeal to complex vectors at all. The construction is formalized in \Lean{}
\cite{Lean4} and machine-checked; the relevant modules carry zero axioms and
zero \texttt{sorry}s.

\paragraph{The thesis, stated plainly.}
The standard reading treats a SIC-POVM as a configuration in complex Hilbert
space, and treats any logical or combinatorial description as a model of that
configuration. We invert the reading. The symmetric informationally-complete
structure is carried, natively and exactly, by the Belnap four-valued lattice
$\BFour = \{\mathrm{N},\mathrm{T},\mathrm{F},\mathrm{B}\}$
\cite{Belnap1977}. The complex vector in $\mathbb{C}^d$ is one
\emph{representation} of that structure, and the difficulty of producing it is
the difficulty of that representation, not of the structure. What is usually
called ``the SIC-POVM existence problem'' splits cleanly into a structural part,
which we discharge, and a representational part, which is Zauner's conjecture.

\paragraph{Why a skeptic should keep reading.}
Two claims in the previous paragraph are the kind that invite dismissal, so we
make them checkable rather than rhetorical. First, ``the structure is carried
by $\BFour$'' is the content of \Cref{thm:unconditional}, whose proof is a
finite computation the reader can replay. Second, ``the representational part
is Zauner's conjecture'' is the content of \Cref{thm:shadow-is-zauner}, a
definitional identity proved by \texttt{rfl}. No step below asks for
assent that a machine has not already granted. The heterodox reading and the
orthodox verification coincide on the same theorems.

%──────────────────────────────────────────────────────────────────────────────
\section{The SIC-POVM condition and the Weyl--Heisenberg group}
%──────────────────────────────────────────────────────────────────────────────

We recall the standard objects, so that the structural claims later have a
fixed target. Let $\omega_d = e^{2\pi i / d}$. The shift and phase operators on
$\mathbb{C}^d$ are $(X v)_k = v_{k-1}$ and $(Z v)_k = \omega_d^{\,k} v_k$, and
the Weyl--Heisenberg displacement operators are
$D_{a,b} = \omega_d^{\,t} X^a Z^b$. The group they generate, $\WH(d)$, acts on
$\mathbb{C}^d$ by the standard projective representation.

\begin{definition}[SIC-POVM, WH-covariant form]
A unit vector $\psi \in \mathbb{C}^d$ is a \emph{fiducial} for a WH-covariant
SIC-POVM if
\[
\|\psi\|^2 = 1
\qquad\text{and}\qquad
|\langle \psi,\ D_{a,b}\,\psi\rangle|^2 = \frac{1}{d+1}
\quad\text{for all } (a,b) \neq (0,0).
\]
Existence of such a $\psi$ in dimension $d$ is \emph{Zauner's conjecture at
$d$}.
\end{definition}

The two conditions are a norm condition and an equiangularity condition.
Everything below concerns where these two conditions actually live.

%──────────────────────────────────────────────────────────────────────────────
\section{The unconditional structural SIC-POVM}
\label{sec:core}
%──────────────────────────────────────────────────────────────────────────────

\subsection{The Belnap lattice and the multilattice fiducial}

Belnap's four-valued logic \cite{Belnap1977} has truth values
$\BFour = \{\mathrm{N},\mathrm{T},\mathrm{F},\mathrm{B}\}$, read as
\emph{neither}, \emph{true}, \emph{false}, and \emph{both}. It carries two
lattice orders (truth and information) and an involution $\neg$ fixing
$\mathrm{N}$ and $\mathrm{B}$. The value $\mathrm{B}$, ``both true and false,''
is the paraconsistent fixed point: $\neg \mathrm{B} = \mathrm{B}$.

For $n \geq 1$ set the state space $\mathrm{MLState}(n) = \BFour^{\,n}$ and the
fiducial $\mathrm{mlFiducial}(n) = \mathrm{B}^{\otimes n}$, the all-$\mathrm{B}$
word. The group $\WH(2)^n$ acts coordinatewise, where the single-register
action of $\WH(2) = (\mathbb{Z}/2)^2$ permutes the four Belnap values. Define
an integer-valued \emph{Frobenius inner product}
$\frobInner$ on $\mathrm{MLState}(n)$ by summing a per-register evidence
functional; write $\langle s, t\rangle_{\mathbb{N}}$ for its value.

\subsection{The theorem}

The following is proved in \Lean{} in the module
\texttt{SIC\_Multilattice\_Proof}. Each conjunct is established by definitional
equality or finite case analysis. The development carries zero axioms and zero
\texttt{sorry}s.

\begin{theorem}[Unconditional structural SIC-POVM]
\label{thm:unconditional}
For every $n \geq 1$, writing $d = 2^n$, the fiducial $\mathrm{B}^{\otimes n}$
satisfies:
\begin{enumerate}
\item \textbf{Cardinality.} The $\WH(2)^n$ orbit has exactly $4^n = d^2$
distinct states.
\item \textbf{Equiangularity.} For every group element $g$,
$\langle \mathrm{B}^{\otimes n},\ g\cdot \mathrm{B}^{\otimes n}\rangle_{\mathbb{N}}
= 2n$, a constant independent of $g$.
\item \textbf{Meet-identity and self-duality.}
$\mathrm{B}^{\otimes n}$ is the identity for the information meet and is fixed
by the involution: $\neg(\mathrm{B}^{\otimes n}) = \mathrm{B}^{\otimes n}$.
\item \textbf{Frobenius closure.} $\frob$ holds registerwise:
$\mathrm{meet}(x,x) = x$ for all states $x$.
\item \textbf{Structural Born ratio.} Every classical (all-$\mathrm{T}$/$\mathrm{F}$)
outcome carries measurement cost $n$, exactly half the fiducial cost $2n$, a
uniform $2\!:\!1$ ratio.
\end{enumerate}
\end{theorem}

The \Lean{} statement is a single conjunction, discharged by a tuple of
finite proofs:

\begin{lstlisting}
theorem sic_povm_belnap_unconditional (n : Nat) :
    (mlOrbit n).card = 4 ^ n
  /\ (forall x, wordMeet (allBWord n) x = x)
  /\ (forall v, (forall i, v i = .T \/ v i = .F) -> totalMeasureCost v = n)
  /\ (forall x, wordJoin (allBWord n) x = allBWord n)
  /\ wordNot (allBWord n) = allBWord n
  /\ (forall x, wordMeet x x = x)
  /\ (forall g h, g <> h -> whAct g (mlFiducial n) <> whAct h (mlFiducial n))
  /\ (forall v, (forall i, v i = .T \/ v i = .F)
        -> mlCost (mlFiducial n) = 2 * mlCost v)
  /\ (forall g, frobInner (mlFiducial n) (whAct g (mlFiducial n)) = 2 * n)
\end{lstlisting}

The three defining features of a SIC-POVM, that the symmetry group produces
$d^2$ frame elements, that all off-diagonal overlaps are equal, and that the
fiducial is self-dual, are present here with none of the apparatus of
$\mathbb{C}^d$. The overlaps are literally equal, as integers; there is no
numerical tolerance and no existence claim to hedge.

%──────────────────────────────────────────────────────────────────────────────
\section{The representational shadow}
\label{sec:shadow}
%──────────────────────────────────────────────────────────────────────────────

\Cref{thm:unconditional} exhibits the structure. What it does not exhibit is a
vector in $\mathbb{C}^d$. We now show that the gap between the two is precisely
Zauner's conjecture, and nothing more.

\subsection{Two inner products at two levels}

The structural theorem and the complex definition speak different inner
products over different objects. This is the whole of the matter.

\begin{center}
\begin{tabularx}{\linewidth}{@{}l X X@{}}
\toprule
 & \textbf{Structural (Belnap)} & \textbf{Representational (complex)} \\
\midrule
State       & $\BFour^{\,n}$ (four-valued word) & $\mathbb{C}^d$ (unit vector) \\
Inner product & $\langle\cdot,\cdot\rangle_{\mathbb{N}}$ (integer-valued) & $\langle\cdot,\cdot\rangle_{\mathbb{C}}$ (Hermitian) \\
Equiangularity & $\langle f, g f\rangle_{\mathbb{N}} = 2n$ (exact) & $|\langle\psi, D\psi\rangle|^2 = \tfrac{1}{d+1}$ \\
Status & proved, zero axioms & Zauner's conjecture \\
\bottomrule
\end{tabularx}
\end{center}

\subsection{The representation problem is Zauner's conjecture}

Let $\mathrm{HilbertEmbeddingExists}(n)$ name the statement that the Belnap
multilattice of \Cref{sec:core} admits a representation as a fiducial vector in
$\mathbb{C}^{2^n}$ carrying the metric equiangularity condition. The following
is proved by definitional identity in the module
\texttt{ZaunerEmbeddingEquivalence} (zero axioms).

\begin{theorem}[The shadow is Zauner]
\label{thm:shadow-is-zauner}
For every $n$, $\mathrm{HilbertEmbeddingExists}(n)$ holds if and only if
Zauner's conjecture holds at $d = 2^n$. The two propositions are definitionally
equal; the equivalence is proved by \texttt{rfl}.
\end{theorem}

\begin{remark}
\Cref{thm:shadow-is-zauner} relocates the open problem. What is open is not
whether the symmetric informationally-complete structure exists; that is
settled by \Cref{thm:unconditional}. What is open is whether that structure
admits a faithful representation in complex Hilbert space with the Hermitian
metric on the nose. The difficulty is a property of the target
$\mathbb{C}^d$, not of the SIC-POVM.
\end{remark}

%──────────────────────────────────────────────────────────────────────────────
\section{One representation route: the mixed-signature Stark conjecture}
%──────────────────────────────────────────────────────────────────────────────

For completeness we record the classical arithmetic-geometry route to the
$\mathbb{C}^d$ vector, and we mark exactly where it borrows. The route is
conditional and its inputs are named; it is one attempt to realize the
representation whose obstruction \Cref{thm:shadow-is-zauner} identified, not an
independent foundation.

Set $m_d = d(d-2)$ and $F_d = \mathbb{Q}(\sqrt{m_d})$, real quadratic for
$d \geq 3$, with ray class field $K_d$. The construction takes a Stark unit
$\varepsilon_d \in K_d^\times$, forms the candidate fiducial from its complex
embeddings, and asks that the resulting vector meet the norm and
equiangularity conditions. In the \Lean{} module \texttt{SIC\_POVM\_Stark} the
field-theoretic objects $F_d, K_d, \mathrm{Gal}(K_d/F_d), \varepsilon_d$ and
the embeddings are declared as typed constants, and the mixed-signature Stark
conjecture is a hypothesis, not an assertion. Two further named axioms carry
the analytic passage from that hypothesis to the metric norm and
equiangularity of the candidate vector. Conditional on the Stark conjecture,
the module proves a fiducial vector exists for all $d \geq 2$.

\begin{theorem}[Conditional representation, after Appleby et al.\ \cite{Appleby2017}]
Assume the mixed-signature Stark conjecture for $K_d$. Then a Weyl--Heisenberg
covariant SIC-POVM vector exists in $\mathbb{C}^d$ for every $d \geq 2$.
\end{theorem}

We stress the honest bookkeeping. The field objects are standard; they are
axiomatized only to avoid importing the full class-field-theory stack. The two
analytic axioms are the genuine open content: they are the point at which the
representation route, working inside $\mathbb{C}^d$, must reconstruct by hand
the equiangularity that \Cref{thm:unconditional} already carries structurally.
That reconstruction is the Zauner residue of \Cref{thm:shadow-is-zauner}. None
of it is load-bearing for \Cref{thm:unconditional}.

%──────────────────────────────────────────────────────────────────────────────
\section{The verification environment: a paraconsistent kernel}
\label{sec:kernel}
%──────────────────────────────────────────────────────────────────────────────

The development is checked under a fork of the Lean 4 (v4.28.0) kernel in which
the principle of explosion is disabled at the C++ level. Standard Lean derives
any goal from a contradiction through the recursor
$\texttt{False.rec} : (C : \mathrm{Sort}\ u) \to \mathrm{False} \to C$; the fork
intercepts recursor and \texttt{casesOn} formation on empty \texttt{Prop}
inductives at three points (\texttt{kernel/type\_checker.cpp},
\texttt{library/constructions/cases\_on.cpp},
\texttt{kernel/environment.\{h,cpp\}}), so that in paraconsistent mode
\texttt{False.elim}, \texttt{False.rec}, \texttt{False.casesOn}, and
\texttt{absurd} are all rejected while every other construct, and \texttt{False}
itself, remain. We flag this prominently because it is the kind of choice a
careful reader should interrogate, and because two facts about it, both
checkable, make it a strengthening of the verification rather than a weakening.

\begin{enumerate}
\item \textbf{Disabling explosion can only remove theorems, never add them.}
Ex falso is an inference rule; deleting it makes strictly fewer propositions
provable and cannot make a false proposition provable. A result that still
type-checks under the stricter kernel is therefore more trustworthy, not less.
A skeptic worried that a nonstandard kernel was tuned to force a result through
has the concern exactly inverted here: this kernel proves less.
\item \textbf{The results of this paper never invoke explosion.}
\Cref{thm:unconditional} and \Cref{thm:shadow-is-zauner} are discharged only by
\texttt{native\_decide}, \texttt{decide}, \texttt{rfl}, and finite case
analysis; none of \texttt{False.elim}, \texttt{False.rec}, or \texttt{absurd}
appears in the SIC modules. The structural results therefore hold with or
without ex falso, and in particular they also compile under the unmodified
Lean 4 kernel.
\end{enumerate}

The kernel is not merely tolerated but exercised. Under the fork it machine-checks
the Belnap identity
\[
\mathrm{band}\ \mathbf{B}\ (\lnot \mathbf{B}) = \mathbf{B} \neq \mathbf{F},
\]
the exact statement a classical kernel collapses by explosion. This is the point
of using it: the SIC fiducial is the value $\mathbf{B}$, the dialetheic fixed
point with $\lnot\mathbf{B} = \mathbf{B}$, and a kernel that trivializes on
contradiction cannot host $\mathbf{B}$ as an inhabited element. The paraconsistent
fork is what lets one and the same environment prove the structural SIC-POVM and
tolerate the fixed point that SIC-POVM is built upon, without triviality. The
structural results occupy the explosion-free fragment common to both kernels.

That the classical kernel is the constrained one, and not the other way round,
admits a precise categorical form, itself machine-checked
(\texttt{ClassicalRestriction.lean}). Take Belnap FOUR in its truth order and
its classical fragment $\{v : v \neq \mathbf{B}\}$. The inclusion of the fragment
is left adjoint to the map $\mathrm{classicalSwitch}$ that collapses
$\mathbf{B} \mapsto \mathbf{F}$, so the classical fragment is a \emph{coreflective}
subcategory of the bilattice, with that collapse as the coreflector and the
identity $\mathrm{classicalSwitch}\circ\text{incl} = \mathrm{id}$ as the unit. The
reverse adjunction fails, and in the information order no adjunction exists at all.
Disabling ex falso is therefore literally a passage to this subcategory, a
corestriction of the paraconsistent ambient, not an extension of it.

%──────────────────────────────────────────────────────────────────────────────
\section{What is proved and what is not}
%──────────────────────────────────────────────────────────────────────────────

We separate the ledger explicitly, because the value of the paper lies in the
separation.

\begin{itemize}
\item \textbf{Proved, unconditionally, zero axioms, zero \texttt{sorry}s:} the
structural SIC-POVM for all $d = 2^n$ (\Cref{thm:unconditional}); the
identification of the complex-representation problem with Zauner's conjecture
(\Cref{thm:shadow-is-zauner}).
\item \textbf{Axiomatized with standard mathematical content:} the field
objects $F_d, K_d$ and their Galois and embedding structure. Replaceable by
existing formalized class field theory; nothing conjectural.
\item \textbf{Open, and honestly marked:} the mixed-signature Stark conjecture,
and with it the analytic passage to the $\mathbb{C}^d$ metric. Equivalently, by
\Cref{thm:shadow-is-zauner}, Zauner's conjecture at $d = 2^n$. This is the
representational residue and nothing in the structural result depends on it.
\end{itemize}

%──────────────────────────────────────────────────────────────────────────────
\section{Discussion}
%──────────────────────────────────────────────────────────────────────────────

The received order of explanation places complex Hilbert space at the
foundation and treats logic as commentary on it. \Cref{thm:unconditional} and
\Cref{thm:shadow-is-zauner} together reverse that order for the symmetric
informationally-complete measurement. The structure that a SIC-POVM is asked to
instantiate is exhibited, exactly and unconditionally, in a four-valued
paraconsistent lattice; the complex vector is one representation of it, and the
long-standing difficulty of producing that vector is the difficulty of the
representation, localized by \Cref{thm:shadow-is-zauner} and not by anything in
the structure.

Two consequences follow, and we state them as consequences of the theorems
above rather than as independent claims. First, bivalence is not where the
ideal quantum measurement lives. The fiducial is the value $\mathrm{B}$, the
point at which true and false coincide; a two-valued logic cannot name it, and
the structure does not appear until the fourth value is available. Second,
complex Hilbert space is a representation target, not a ground. It is the arena
in which the reconstruction is hard, not the arena in which the fact is true.

A reader inclined to read foundations widely may take a further step, which we
mark as interpretation rather than theorem: if the optimally informative
measurement, the very act by which empirical inquiry extracts the most it can
from a system, is grounded in a four-valued paraconsistent logic rather than in
the continuum it is usually phrased over, then the priority of the continuum
and of bivalence in our account of the physical world is a representational
habit and not a necessity. The measurement was always four-valued. We have
only now been able to check it by machine.

%──────────────────────────────────────────────────────────────────────────────
\section*{Data and code availability}
%──────────────────────────────────────────────────────────────────────────────

The \Lean{} development (modules \texttt{SIC\_Multilattice\_Proof},
\texttt{BelnapNFiducial}, \texttt{QCI\_SICPOVM\_Bridge},
\texttt{ZaunerEmbeddingEquivalence}, \texttt{SIC\_POVM\_ParityGate},
\texttt{SIC\_POVM\_Functor}, \texttt{SIC\_POVM\_Stark}) builds under a single
\texttt{lake build} and is available in the \texttt{p4rakernel} repository
\cite{p4rakernel}, whose \texttt{src/} is the paraconsistent kernel fork of
\Cref{sec:kernel} (built via \texttt{make}, then \texttt{lake build} against the
resulting binary). \Cref{thm:unconditional} and \Cref{thm:shadow-is-zauner}
reside in modules carrying zero axioms and zero \texttt{sorry}s, and use only
explosion-free tactics. We have checked the SIC modules directly under both
kernels: they type-check unchanged under the forked binary
(\texttt{build/stage1/bin}) and under stock Lean 4.

%──────────────────────────────────────────────────────────────────────────────
\section*{Acknowledgements}
%──────────────────────────────────────────────────────────────────────────────

The author would like to thank Harry T.\ Larson, for imparting the importance
of catching rising problems, and never letting them go \cite{Larson1961}.

%──────────────────────────────────────────────────────────────────────────────
\begin{thebibliography}{99}

\bibitem{Appleby2013}
D.\ M.\ Appleby, I.\ Bengtsson, S.\ Flammia, and D.\ Grassl,
\textit{The Monomial Representations of the Clifford Group},
Quantum Inf.\ Comput.\ \textbf{12} (2012), 404--431.

\bibitem{Appleby2017}
D.\ M.\ Appleby, T.-Y.\ Chien, S.\ Flammia, and S.\ Waldron,
\textit{Constructing Exact Symmetric Informationally Complete Measurements from
Numerical Solutions},
J.\ Phys.\ A: Math.\ Theor.\ \textbf{51} (2018), 165302.

\bibitem{Belnap1977}
N.\ D.\ Belnap,
\textit{A Useful Four-Valued Logic},
in: J.\ M.\ Dunn and G.\ Epstein (eds.), \textit{Modern Uses of Multiple-Valued Logic},
Reidel, Dordrecht, 1977, pp.\ 5--37.

\bibitem{Larson1961}
H.\ T.\ Larson,
\textit{Catch a Rising Problem\ldots\ and Never Ever Let it Go},
IRE Transactions on Engineering Management \textbf{EM-8} (1961), no.\ 4, 173--174.
\url{https://ieeexplore.ieee.org/document/1641382}

\bibitem{Lean4}
L.\ de Moura and S.\ Ullrich,
\textit{The Lean 4 Theorem Prover and Programming Language},
in: CADE-28, Lecture Notes in Comput.\ Sci.\ vol.\ 12699, Springer, 2021, pp.\ 625--635.

\bibitem{p4rakernel}
C.\ L.\ Mills,
\textit{p4rakernel: Lean\,4 paraconsistent kernel and formalization},
software repository, 2026.

\bibitem{Zauner1999}
G.\ Zauner,
\textit{Quantum designs: Foundations of a noncommutative design theory},
PhD thesis, University of Vienna, 1999.

\end{thebibliography}

\end{document}
